\section{Conclusion}
\label{sec:conclusion}

If income moves from labour to capital, the tax system has to collect on the second what it loses
on the first. We stated that requirement as a condition, assembled both of its sides from sourced
components, and measured how far the existing system falls short.

On a consistently federal basis the effective marginal rate on income shifting from wages to
capital is \result{TauKFederalOnly} percent, against the \result{RequiredFedEasy} to
\result{RequiredFedHard} percent replacement requires. Written with an explicit pass-through
coefficient, the comparison becomes a boundary, and the boundary is the more useful object: the
condition closes only where taxable capital income per displaced wage dollar reaches
\result{GStarFedEasy} to \result{GStarFedHard}, above the ceiling of one that output-preserving
displacement implies. Making the comparison consistent across jurisdictions widens the gap rather
than closing it, because the state adds more to the labour side than to the capital side.

The reason the rate is low is the contribution. It is not profit shifting, which ranks third of
seven levers and moves the assembled rate by \result{SwingShifted} across its entire range,
against \result{SwingDeferral} for deferral and step-up at death. What binds is ownership. About
\result{ThetaUntaxable} percent of US corporate equity sits in forms the individual income tax
barely reaches, and of the gains that do accrue in taxable accounts,
\result{GainsUntaxedAtDeath} percent are never taxed at all. A legislature can change a rate. It
cannot easily change who owns the claims, and the rate operates on a base most households are
not in.

That single fact does more work than one result. The same ownership concentration means
household debt-service capacity cannot be restored by broadening ownership, since
\result{HHEquityOutsideHouseholds} percent of corporate equity sits outside the household sector
and broadened ownership without liquidity moves the boundary by exactly zero. It means an equity
bust is a large fiscal event and a small credit event, because the assets that fall are held
where the propensity to consume is lowest. The revenue problem, the incidence problem and the
transmission result are three readings of one structural fact.

We do not claim this settles which tax base is right, and we state the case against our own
choice where it arises. Three components are swept rather than sourced, the pass-through
coefficient cannot be sourced at all, and the household figures are incidence arithmetic rather
than forecasts. What holds on every base and under every reading we can defend is that the
central case falls short, that the shortfall is governed by ownership rather than by enforcement,
and that the instruments which would actually move it are the ones that have had the least
attention.
